\chapter{Conclusão}

O ponto de partida foi um problema concreto: as informações financeiras de uma pessoa estão espalhadas por vários aplicativos, e transformá-las em decisão depende da iniciativa e da interpretação do próprio usuário. A pesquisa realizada mostrou que 83\% dos respondentes mantinham duas ou mais contas, que metade não conhecia o Open Finance e que a preocupação com a privacidade era alta. O EIGAM foi desenvolvido para responder a esse cenário, e o resultado é um assistente financeiro em funcionamento, que consolida os dados de várias instituições pelo Open Finance e os entrega em linguagem natural no mensageiro que a pessoa já usa.

O principal ponto forte do produto diante desse problema é reduzir o esforço do usuário sem reduzir o seu controle. A pergunta ``quanto gastei no mercado este mês, somando todos os bancos?'' deixa de exigir a abertura de três aplicativos e um cálculo manual: ela é respondida em uma única mensagem, com o valor calculado em centavos sobre dados efetivamente consultados, e não estimado pelo modelo de linguagem. A proatividade prometida está presente nos resumos periódicos, entregues na faixa de horário escolhida pelo titular, e implementada nos avisos de limite de gasto, de risco de ultrapassagem do limite, de cobrança duplicada, de transação atípica e de aumento por categoria, com entrega durável e com autorização reconferida no instante do envio. A mesma lógica de reduzir esforço levou a conversa até onde o usuário já está: além do WhatsApp, do Telegram e do painel \textit{web}, o titular pode conectar o EIGAM ao assistente de inteligência artificial que já usa, por um servidor MCP somente de leitura que obedece às mesmas regras de máscara e auditoria.

O segundo ponto forte é a confiança sustentada por mecanismos técnicos verificáveis. Os valores ficam mascarados até o PIN, por uma regra do servidor que o modelo não consegue contornar; os comandos da LGPD são reconhecidos por regras determinísticas, com 100\% de revocação medida; os dados sensíveis são cifrados campo a campo; e cada direito do titular pode ser exercido com uma frase, inclusive a exclusão completa da conta. Quando os dados não sustentam uma resposta, o sistema se abstém em vez de exibir um número com aparência de precisão: as projeções são apresentadas apenas pelo método aprovado no teste retrospectivo e identificadas como experimentais, e o saldo livre abstém-se, com motivo declarado, quando faltam histórico ou dados atualizados. Como, na pesquisa, menos da metade dos respondentes (47,6\%) declarou confiança alta em assistentes de inteligência artificial para tratar de finanças, essa transparência é parte da proposta de valor.

O terceiro ponto forte é a arquitetura. A decisão de desacoplar o núcleo do canal de mensageria, defendida na fase de projeto como forma de reduzir a dependência de um único fornecedor, foi comprovada na prática com a inclusão do Telegram e da conversa no painel \textit{web} sem alterações nos serviços de domínio, e com o servidor MCP, que reaproveita as mesmas rotas de leitura do painel. A modularidade da arquitetura também permitiu ajustar a modelagem à medida que a medição trazia evidências, como na substituição do Prophet e do classificador de intenções e no aproveitamento da categoria detalhada fornecida pelo próprio agregador de Open Finance, sem reescrever o sistema. O teste de carga estimou que um único servidor de 4 vCPU atende cerca de 2.850 usuários sem formação de fila, ao custo de cerca de 12 euros mensais de infraestrutura.

Os objetivos específicos apresentados no Capítulo~\ref{chap:contexto} foram atendidos. As contas bancárias do titular são integradas pelo Open Finance, mediante consentimento, com várias instituições por titular e revogação pela própria conversa. As consultas em linguagem natural têm os valores calculados por buscas estruturadas sobre os dados sincronizados, sem que o modelo de linguagem gere números. Os resumos periódicos são entregues de forma proativa na faixa de horário escolhida pelo titular, e os avisos e as projeções estão implementados, ainda sem emissão ao público geral. A proteção dos dados combina o mascaramento por PIN, a cifragem campo a campo e o exercício, por mensagem, dos direitos do titular previstos na LGPD. Por fim, a qualidade foi medida nas quatro frentes previstas: testes automatizados, conjuntos de referência, teste de carga e teste retrospectivo.

Conclui-se que o objetivo geral foi atingido: foi desenvolvida uma solução de assistência financeira inteligente, preditiva e conversacional, integrada ao Open Finance, que fornece consultas, análises, previsões, alertas e orientações personalizadas em diferentes ferramentas de conversa, a partir de uma arquitetura desacoplada do provedor de mensageria e verificada em produção com múltiplos canais simultâneos. As consultas, as análises, as orientações, as metas e, como estimativa experimental, a projeção do gasto e a simulação de compra estão em uso no beta; os alertas estão implementados, com a liberação ao público geral condicionada a etapas de validação. A privacidade do titular foi tratada como requisito de projeto desde a concepção.

Essas etapas e as demais limitações, discutidas na Seção~\ref{sec:limitacoes}, não invalidam a proposta, mas orientam os trabalhos futuros, apresentados no capítulo seguinte.
